\sect{पुरुषोत्तम-मासस्य कृष्ण-परमा-एकादशी-माहात्म्यम्}
\label{sec:padma-purushottama-krishna-parama}

\ganapatyadiDhyanam
\padmaPuranam

\dnsub{अथ पञ्चषष्टितमोऽध्यायः॥६५}

\uvacha{युधिष्ठिर उवाच}

\twolineshloka
{श्रुतानि बहुधर्माणि व्रतानि च जगत् प्रभो}
{एकादशीसमं किञ्जिच्छ्रुतं नैव जनार्दन}% ॥१॥

\twolineshloka
{पुनस्त्वेकादशीं ब्रूहि पापघ्नीं पुण्यदायिनीम्}
{यां कृत्वा मनुजो लोके प्राप्नुयात् परमं पदम्}% ॥२॥

\uvacha{श्रीकृष्ण उवाच}

\twolineshloka
{शुक्ले वा यदि वा कृष्णे यदा चैकादशी भवेत्}
{न त्याज्या जगतीपाल मोक्षसौख्यविवर्धनी}% ॥३॥

\twolineshloka
{एकादशी कलौ राजन् भवबन्धविमोचनी}
{कामदा सर्वकामानां पापानां पापहा भुवि}% ॥४॥

\twolineshloka
{रविवारेऽथ माङ्गल्ये सङ्क्रमे वा नृपोत्तम}
{एकादशी सदोपोष्या पुत्रपौत्रविवर्धनी}% ॥५॥

\twolineshloka
{एकादशीव्रतं क्वापि न त्याज्यं विष्णुवल्लभैः}
{आयुः कीर्तिप्रदं नित्यं सन्तानारोग्यवित्तदम्}% ॥६॥

\twolineshloka
{मोक्षदं रूपदं राज्यं नित्यमेकादशीव्रतम्}
{ये कुर्वन्ति महीपाल श्रद्धया परया युताः}% ॥७॥

\twolineshloka
{यथोक्तविधिना लोके ते नरा विष्णुरूपिणः }
{जीवन्मुक्तास्तु भूपाल दृश्यन्ते नात्र संशयः}% ॥८॥

\uvacha{युधिष्ठिर उवाच}

\twolineshloka
{जीवन्मुक्ताः कथं कृष्ण विष्णुरूपाः कथं पुनः}
{पापरूपाश्च दृश्यन्ते परं कौतूहलं हि मे}% ॥९॥

\uvacha{श्रीकृष्ण उवाच}

\twolineshloka
{ये च राजन् कलौ भक्त्या निर्जलं व्रतमुत्तमम्}
{एकादश्याः प्रकुर्वन्ति विधिदृष्टेन कर्मणा}% ॥१०॥

\twolineshloka
{न कथं विष्णुरूपास्ते जीवन्मुक्ताः कथं नहि}
{सर्वपापहरं पुण्यं व्रतमेकादशीसमम्}% ॥११॥

\twolineshloka
{न किञ्चिद् विद्यते राजन् सर्वकामप्रदं नृणाम्}
{एकाशनं दशम्यां च नन्दायां निर्जलं व्रतम्}% ॥१२॥

\twolineshloka
{पारणं चैव भद्रायां कृत्वा विष्णुसमा नराः}
{श्रद्धावान्यस्तु कुरुते कामदाया व्रतं शुभम्}% ॥१३॥

\twolineshloka
{वाञ्छितं लभते सोऽपि इहलोके परत्र च}
{पवित्रा पावनी ह्येषा महापातकनाशिनी}% ॥१४॥

\twolineshloka
{भुक्तिमुक्तिप्रदा चैव कर्तॄणां नृपसत्तम}
{कामदायां विधानेन पूजयेत् पुरुषोत्तमम्}% ॥१५॥

\twolineshloka
{पुष्पधूपादिभिश्चैव नैवेद्यैर्विविधैस्तथा}
{कांस्य मांसमसूरांश्च चणकान्कोद्रवांस्तथा}% ॥१६॥

\twolineshloka
{शाकं मधु परान्नं च पुनर्भोजन मैथुनम्}
{वैष्णवो व्रतकर्ता च दशम्यां दश वर्जयेत्}% ॥१७॥

\twolineshloka
{द्यूतं क्रीडां तथा निद्रा ताम्बूलं दन्तधावनम्}
{परापवादं पैशुन्यं स्तेयं हिंसां तथा रतिम्}% ॥१८॥

\twolineshloka
{क्रोधं च वितथं वाक्यमेकादश्यां विवर्जयेत्}
{कांस्यं मांसं मसूरांश्च तैलं वितथभाषणम्}% ॥१९॥

\twolineshloka
{व्यायामं च प्रवासं च पुनर्भोजनमैथुनम्}
{वृषपृष्ठं परान्नं च शाकं च द्वादशीदिने}% ॥२०॥

\twolineshloka
{अनेन विधिना राजन् विहिता यैश्च कामदा}
{रात्रौ जागरणं कृत्वा पूजितः पुरुषोत्तमः}% ॥२१॥

\twolineshloka
{सर्वपापविनिर्मुक्तास्ते यान्ति परमां गतिम्}
{पठनाच्छ्रवणाद् राजन् गोसहस्रफलं लभेत्}% ॥२२॥

॥इति श्रीपाद्मे महापुराणे पञ्चपञ्चाशत्साहस्र्यां संहितायामुत्तरखण्डे उमापतिनारदसंवादान्तर्गतकृष्णयुधिष्ठिरसंवादे पुरुषोत्तम-मासस्य शुक्ल-कामदा-एकादशी-माहात्म्यं नाम पञ्चषष्टितमोऽध्यायः॥६५॥

\genMangalaShloka

\hyperref[sec:ekadashi_mahatmyam_padma_puranam]{\closesub}